% =====================================================================
%  3. LA LÓGICA MODAL K
% =====================================================================
\chapter{La lógica modal \texorpdfstring{$\K$}{K}}\label{sec:logicamodal}

Este capítulo contiene la parte matemática del trabajo en su versión ``de
papel'': la sintaxis del lenguaje modal básico, el sistema axiomático $\K$,
la semántica de Kripke y los enunciados de corrección y completitud, junto
con el mapa de la demostración de esta última. Todo ello sigue en lo
esencial a Blackburn, de Rijke y Venema \cite{blackburn}, con las
adaptaciones, que iremos señalando, que conviene hacer pensando en la
formalización del capítulo \ref{sec:formalizacion}.

\section{¿Necesidad y posibilidad?}

La lógica proposicional clásica trata cada enunciado como verdadero o
falso, sin matices. Pero el lenguaje natural distingue \emph{modos} de
verdad: no es lo mismo ``llueve'' que ``necesariamente llueve'' o que
``es posible que llueva''. La lógica modal introduce dos operadores para
capturar esa distinción: $\nec\varphi$ (``$\varphi$ es necesaria'') y
$\pos\varphi$ (``$\varphi$ es posible''). Ambos están ligados por una
dualidad esperable: algo es posible exactamente cuando su negación no es
necesaria, $\pos\varphi \equiv \lnot\nec\lnot\varphi$. En nuestro ejemplo de lluvia, ``es posible que llueva'' es lo mismo que ``no es necesario que no llueva''.

Parte del atractivo de la lógica modal es que la lectura de los operadores
no está fijada. Según la aplicación, un mismo par $\nec$, $\pos$ admite
interpretaciones muy distintas:

\begin{itemize}
  \item \textbf{Modalidad alética}, la lectura clásica: $\nec\varphi$ es
        ``$\varphi$ es necesariamente cierta'' y $\pos\varphi$, ``$\varphi$
        es posible''.
  \item \textbf{Lógica epistémica}: $\nec\varphi$ es ``se sabe que
        $\varphi$'' y $\pos\varphi$, ``$\varphi$ es compatible con lo que se
        sabe'', es decir, no puede descartarse. Los operadores suelen
        escribirse entonces $K\varphi$ y $\widehat{K}\varphi$.
  \item \textbf{Lógica temporal}: $\nec\varphi$ es ``$\varphi$ será siempre
        cierta en el futuro'' y $\pos\varphi$, ``$\varphi$ será cierta en
        algún momento''.
  \item \textbf{Lógica deóntica}: $\nec\varphi$ es ``$\varphi$ es
        obligatoria'' y $\pos\varphi$, ``$\varphi$ está permitida''.
\end{itemize}

En todos los casos la dualidad se mantiene, y resulta iluminadora: lo
permitido es lo que no está prohibido, lo compatible con el conocimiento es
lo que no se sabe falso, lo que ocurrirá alguna vez es lo que no siempre
será falso. Cada lectura sugiere, en cambio, principios distintos: por
ejemplo, ``lo sabido es cierto'' ($\nec\varphi \limp \varphi$) es razonable
para el conocimiento, y también para el tiempo si se admite el instante
presente, pero no para la obligación (que algo sea obligatorio no lo hace
cierto). Esta multiplicidad explica que exista toda una familia de lógicas
modales.

La que estudiamos aquí, la lógica $\K$ (por Kripke), es la base
común y más débil de todas ellas: no compromete ninguna lectura y sirve de
punto de partida para las demás, que se obtienen añadiendo axiomas. La
lectura epistémica reaparecerá en el capítulo \ref{sec:muddy} AAA.

\section{Sintaxis}\label{sec:sintaxis-mat}

\begin{definicion}[Lenguaje modal básico]\label{def:lenguaje}
Fijado un conjunto $\mathrm{VP}$ de \emph{variables proposicionales}, las
\emph{fórmulas modales} se definen por la siguiente gramática:
\[
  \varphi ::= p \mid \bot \mid \varphi \limp \varphi \mid \nec\varphi,
  \qquad p \in \mathrm{VP}.
\]
Es decir, el conjunto $\mathrm{Form}$ de fórmulas es el menor conjunto que contiene
las variables y la constante $\bot$ (``falso'') y está cerrado por la
implicación y por el operador $\nec$.
\end{definicion}

Adoptamos deliberadamente un lenguaje primitivo mínimo: solo $\bot$,
$\limp$ y $\nec$. El resto de conectivas se introducen como
abreviaturas:
\[
\begin{array}{rclcrcl}
  \top &:=& \bot \limp \bot, &\qquad&
  \lnot\varphi &:=& \varphi \limp \bot,\\[2pt]
  \varphi \lor \psi &:=& \lnot\varphi \limp \psi, &&
  \varphi \land \psi &:=& \lnot(\varphi \limp \lnot\psi),\\[2pt]
  \varphi \leftrightarrow \psi &:=& (\varphi\limp\psi)\land(\psi\limp\varphi), &&
  \pos\varphi &:=& \lnot\nec\lnot\varphi.
\end{array}
\]
Las definiciones de $\lor$ y $\land$ son las clásicas: $\varphi \lor \psi$
dice ``si no $\varphi$, entonces $\psi$'', y $\varphi \land \psi$ niega que
$\varphi$ implique la negación de $\psi$. La ventaja de trabajar con un
lenguaje mínimo es doble. En el plano matemático, las definiciones y
demostraciones por inducción estructural tienen solo cuatro casos. En el
plano de la formalización, cada constructor del tipo de fórmulas es un caso
en cada demostración: menos primitivas significa menos casos que demostrar,
y las conectivas derivadas reciben su significado ``gratis'', como
comprobaremos en la sección \ref{sec:form-semantica}.

Esta es la primera de las adaptaciones. Blackburn, de Rijke y
Venema \cite[Def.~1.9]{blackburn} toman como primitivos $\bot$, $\lnot$,
$\lor$ y $\pos$, y es $\nec$ el que se define, como
$\nec\varphi := \lnot\pos\lnot\varphi$; la implicación es allí la
abreviatura $\varphi \limp \psi := \lnot\varphi \lor \psi$. Nosotros
invertimos la elección: $\nec$ es primitivo y $\pos$ es la abreviatura.
Ambos lenguajes son intertraducibles (cada fórmula de uno se reescribe en
el otro sin alterar su significado respecto a la semántica de la sección
\ref{sec:kripke-mat}), y los propios autores reconocen que su preferencia
por $\pos$ obedece a la comodidad para el trabajo algebraico de capítulos
posteriores del libro. Pensando en la formalización, $\nec$ es la opción
natural: su cláusula semántica es una cuantificación universal, más cómoda
de manejar en un asistente de demostración que la existencial de $\pos$
(sección \ref{sec:form-semantica}). La misma elección hace innecesario uno
de los axiomas del sistema de Blackburn, como veremos en la sección
\ref{sec:sistemaK}.

Obsérvese también que no imponemos ninguna condición sobre el conjunto
$\mathrm{VP}$. Blackburn \cite[p.~10]{blackburn} lo supone habitualmente
numerable, hipótesis que interviene en su demostración del lema de
Lindenbaum; nosotros podremos prescindir de ella (sección
\ref{sec:mapa-completitud}, paso 3).

Como es habitual, la inducción (y la recursión) estructural sobre
fórmulas será nuestra herramienta básica: para probar que toda fórmula tiene
una propiedad basta verla para $p$ y $\bot$, y verla para
$\varphi \limp \psi$ y $\nec\varphi$ supuesta para las subfórmulas.

\section{El sistema axiomático \texorpdfstring{$\K$}{K}}\label{sec:sistemaK}

Un sistema axiomático determina qué
fórmulas son demostrables a partir de unos axiomas mediante unas
reglas de inferencia. El sistema $\K$ toma como punto de partida todo el
cálculo proposicional clásico y le añade lo mínimo para razonar con $\nec$.

\begin{definicion}[El sistema $\K$]\label{def:sistemaK}
La relación de demostrabilidad $\proves \varphi$ (``$\varphi$ es un teorema
de $\K$'') es la menor relación que cumple:
\begin{itemize}
  \item[\textup{(CPL)}] $\proves \varphi$ para toda \emph{tautología
        clásica} $\varphi$ (véase más abajo);
  \item[\textup{(K)}] $\proves \nec(\varphi \limp \psi) \limp
        (\nec\varphi \limp \nec\psi)$, para cualesquiera fórmulas
        $\varphi, \psi$ (\emph{axioma de distribución});
  \item[\textup{(MP)}] si $\proves \varphi \limp \psi$ y $\proves \varphi$,
        entonces $\proves \psi$ (\emph{modus ponens});
  \item[\textup{(N)}] si $\proves \varphi$, entonces $\proves \nec\varphi$
        (\emph{regla de necesitación}).
\end{itemize}
\end{definicion}

Detengámonos en (CPL). Lo habitual en los textos es dar tres esquemas de
axiomas proposicionales concretos (los axiomas de Hilbert) y obtener las
tautologías como teoremas. Aquí seguimos el camino, también estándar
\cite[p.~33]{blackburn}, de admitir directamente todos los esquemas
de tautologías; los propios autores lo describen como un ``exceso
axiomático'' (\emph{axiomatic overkill}) que prefieren a la alternativa de
fijar unos pocos axiomas y derivar de ellos el resto, refinamiento que, a
su juicio, tiene poco interés para sus propósitos, y lo mismo vale para los
nuestros. Precisemos la noción: una fórmula modal es una \emph{tautología
clásica} si es verdadera bajo toda valoración booleana que asigne valores
de verdad arbitrarios a sus variables y a sus subfórmulas maximales
de la forma $\nec\psi$, tratadas estas como si fueran átomos opacos. La
valoración no mira dentro de $\nec\psi$: en $\nec\nec p \limp \nec p$, por
ejemplo, los ``átomos'' son $\nec\nec p$ y $\nec p$, que reciben valores
independientes, y por eso la fórmula no es una tautología clásica. Es fácil
convencerse de que esta definición equivale a la de Blackburn, ``toda
instancia de una tautología proposicional'': basta sustituir cada
subfórmula maximal de la forma $\nec\psi$ por una variable nueva (la misma
para todas las ocurrencias de la misma subfórmula). Así,
$\nec\psi \limp \nec\psi$ es una tautología clásica (es una instancia de
$\alpha \limp \alpha$), pero $\nec p \limp p$, con $p$ una variable, no lo
es: hay valoraciones que hacen verdadero el ``átomo'' $\nec p$ y falsa
$p$.\footnote{Conviene escribir $\nec p \limp p$ y no
$\nec\varphi \limp \varphi$: si $\varphi$ es ella misma una tautología,
como $\top$, entonces $\nec\varphi \limp \varphi$ sí es una tautología
clásica, porque ninguna valoración hace falsa $\varphi$.} Esta
presentación evita desarrollar dentro de $\K$ toda la
maquinaria del cálculo proposicional, que no es el objeto del trabajo, y
tiene una traducción muy limpia a Lean, donde ``ser tautología'' se puede
definir semánticamente (sección \ref{sec:form-tautologias}). El
precio es que (CPL) es un esquema infinito, algo que a un sistema formal
como Lean no le supone ningún problema.

El sistema de la definición \ref{def:sistemaK} no es literalmente el $\K$
de Blackburn \cite[Def.~1.39]{blackburn}, y conviene dejar claro en qué
difieren y por qué tienen los mismos teoremas. Allí $\K$ cuenta con un
tercer axioma, \textup{(Dual)} $\pos p \leftrightarrow \lnot\nec\lnot p$, y
con una tercera regla, la \emph{sustitución uniforme}: de
$\proves \varphi$ se pasa a $\proves \varphi^\sigma$, donde
$\varphi^\sigma$ resulta de sustituir uniformemente las variables de
$\varphi$ por fórmulas cualesquiera. Además, los axiomas (K) y (Dual) se
enuncian con variables concretas $p$, $q$, y no como esquemas.
\begin{itemize}
  \item \textup{(Dual)} sobra en nuestro lenguaje. Como aquí $\pos\varphi$
        es por definición $\lnot\nec\lnot\varphi$, el axioma se
        desplegaría a $\lnot\nec\lnot\varphi \leftrightarrow
        \lnot\nec\lnot\varphi$, una instancia de
        $\alpha \leftrightarrow \alpha$, que ya está en (CPL). Los propios
        autores lo advierten \cite[p.~34]{blackburn}: (Dual) solo hace
        falta porque su primitivo es $\pos$, y de haber tomado $\nec$ como
        primitivo no habría sido necesario.
  \item La \emph{sustitución uniforme} sobra porque nuestros axiomas son
        esquemas: (K) se postula de golpe para todas las fórmulas
        $\varphi, \psi$, y el conjunto de las tautologías es cerrado bajo
        sustitución por su propia definición. Con ello, la cerradura de los
        teoremas bajo sustitución deja de ser una regla y pasa a ser un
        teorema, que se demuestra por inducción sobre la derivación:
        las instancias de tautologías son tautologías, las instancias de
        (K) son instancias de (K), y las reglas (MP) y (N) conmutan con la
        sustitución.
\end{itemize}
Así pues, formulados en nuestro lenguaje, el sistema de Blackburn y el de
la definición \ref{def:sistemaK} demuestran exactamente las mismas
fórmulas: cada uno contiene los axiomas del otro y es cerrado bajo sus
reglas. Dicho de otro modo, el conjunto de nuestros teoremas contiene las
tautologías, (K) y (Dual), y es cerrado bajo (MP), (N) y sustitución
uniforme, es decir, es una \emph{lógica modal normal} en el sentido de
\cite[Def.~1.42]{blackburn}; y, al estar generado por lo imprescindible, es
la menor de todas ellas, que es justamente como se define $\K$ en esa
definición. De hecho, nuestra definición ``como la menor relación que
cumple\dots'' está más cerca de esa caracterización conjuntista que de la
definición 1.39 (demostraciones como sucesiones finitas de fórmulas); la
equivalencia entre ambos puntos de vista es el ejercicio 1.6.6 de
\cite{blackburn}, y la versión ``menor relación'' es precisamente la que
produce una definición inductiva en Lean.

Los axiomas (K) y la regla (N) capturan el comportamiento mínimo de $\nec$:
la necesidad distribuye sobre la implicación, y los teoremas son necesarios.
Nótese lo que la regla (N) no dice: no afirma
$\proves \varphi \limp \nec\varphi$ (las verdades contingentes no son
necesarias), sino solo que las fórmulas demostrables pueden
``encajonarse''. Blackburn dedica a este punto una advertencia expresa
\cite[Obs.~1.41]{blackburn}: quien venga de la deducción natural puede
sentirse tentado de suponer $p$, aplicar (N) para obtener $\nec p$ y
``descargar'' la hipótesis concluyendo $p \limp \nec p$, que no es válida.
El error está en aplicar (N) a una hipótesis: la regla solo se aplica a
teoremas. Volveremos sobre ello al definir la deducción desde hipótesis
(sección \ref{sec:mapa-completitud}).

De estos ingredientes se derivan reglas útiles. Destacamos las dos que
formalizaremos, con sus demostraciones en papel, que después podrán
compararse línea a línea con el código:

\begin{proposicion}[Regla de monotonía, RM]\label{prop:rm}
Si\/ $\proves \varphi \limp \psi$, entonces
$\proves \nec\varphi \limp \nec\psi$.
\end{proposicion}
\begin{proof}
Por (N), de $\proves \varphi \limp \psi$ obtenemos
$\proves \nec(\varphi \limp \psi)$. El axioma (K) da
$\proves \nec(\varphi \limp \psi) \limp (\nec\varphi \limp \nec\psi)$, y
(MP) concluye.
\end{proof}

\begin{proposicion}[Regla de equivalencia, RE]\label{prop:re}
Si\/ $\proves \varphi \leftrightarrow \psi$, entonces
$\proves \nec\varphi \leftrightarrow \nec\psi$.
\end{proposicion}
\begin{proof}
De $\proves \varphi \leftrightarrow \psi$ se extraen, usando las tautologías
$(\varphi \leftrightarrow \psi) \limp (\varphi \limp \psi)$ y
$(\varphi \leftrightarrow \psi) \limp (\psi \limp \varphi)$ junto con (MP),
las dos implicaciones $\proves \varphi \limp \psi$ y
$\proves \psi \limp \varphi$. La regla RM da entonces
$\proves \nec\varphi \limp \nec\psi$ y $\proves \nec\psi \limp \nec\varphi$,
y la tautología $(\alpha \limp \beta) \limp ((\beta \limp \alpha) \limp
(\alpha \leftrightarrow \beta))$, con dos aplicaciones de (MP), las reúne en
la equivalencia buscada.
\end{proof}

Obsérvese el patrón de estas demostraciones: los únicos movimientos
permitidos son citar un axioma y aplicar una regla. Es un estilo rígido pero es exactamente esta rigidez
la que lo hace formalizable de manera directa.

\section{Semántica de Kripke}\label{sec:kripke-mat}

La semántica estándar de la lógica modal, debida a Kripke \cite{kripke}, se apoya en una
idea muy visual: la de \emph{mundos posibles} conectados por una relación de
accesibilidad. Una fórmula no es verdadera o falsa en abstracto, sino
en un mundo; y los operadores modales cuantifican sobre los mundos
accesibles desde el actual.

\begin{definicion}[Marcos y modelos]\label{def:kripke}
Un \emph{marco de Kripke} es un par $F = (W, R)$ formado por un conjunto no
vacío\footnote{En la formalización el tipo de mundos podrá ser vacío; véase
la observación \ref{obs:vacio} en la sección \ref{sec:form-semantica}.} $W$
de mundos y una relación binaria $R \subseteq W \times W$, la
\emph{relación de accesibilidad}. Un \emph{modelo de Kripke} sobre
$\mathrm{VP}$ es un par $M = (F, V)$ donde $F$ es un marco y
$V \colon W \to \mathcal{P}(\mathrm{VP})$ es una \emph{valoración} que
indica qué variables son verdaderas en cada mundo.
\end{definicion}

Un marco es, sencillamente, un grafo dirigido; escribimos $R\,w\,v$ o
$w \mathrel{R} v$ para ``desde $w$ es accesible $v$''. En la valoración
hay otra pequeña adaptación respecto a \cite[Def.~1.19]{blackburn}, donde
$V \colon \mathrm{VP} \to \mathcal{P}(W)$ asigna a cada variable el
conjunto de mundos en que es verdadera, de modo que la cláusula atómica de
la definición siguiente se lee allí $w \in V(p)$. Nosotros la escribimos al
revés, $V \colon W \to \mathcal{P}(\mathrm{VP})$, y leeremos $p \in V(w)$.
La información es la misma (en ambos casos, un subconjunto de
$W \times \mathrm{VP}$), pero nuestra dirección es la que usaremos en la
figura \ref{fig:kripke-ejemplo}, en el modelo canónico de la sección
\ref{sec:mapa-completitud} (donde la valoración de un mundo $\Gamma$ es el
conjunto de variables que pertenecen a $\Gamma$) y en el código, donde la
valoración recibe primero el mundo y después la variable.

\begin{definicion}[Satisfacción]\label{def:satisfaccion}
Dado un modelo $M$ y un mundo $w$, la relación $M, w \vDash \varphi$
(``$\varphi$ es verdadera en el mundo $w$ de $M$'') se define por recursión
sobre $\varphi$:
\[
\begin{array}{lcl}
  M, w \vDash p &\iff& p \in V(w),\\
  M, w \vDash \bot && \text{no se da nunca},\\
  M, w \vDash \varphi \limp \psi &\iff& \text{si } M, w \vDash \varphi
      \text{ entonces } M, w \vDash \psi,\\
  M, w \vDash \nec\varphi &\iff& \text{para todo } v \text{ con } R\,w\,v,\
      M, v \vDash \varphi.
\end{array}
\]
\end{definicion}

La cláusula interesante es la última: $\nec\varphi$ es verdadera en $w$
cuando $\varphi$ lo es en todos los mundos accesibles desde $w$.
Desplegando las abreviaturas se comprueba que las conectivas derivadas
reciben el significado esperado; en particular,
\[
  M, w \vDash \pos\varphi \iff \text{existe } v \text{ con } R\,w\,v
  \text{ y } M, v \vDash \varphi:
\]
la posibilidad es la cuantificación existencial sobre los mundos accesibles.

\begin{figure}[htb]
\centering
\begin{tikzpicture}[
    world/.style={circle, draw, minimum size=1.05cm, inner sep=1pt},
    flecha/.style={-{Stealth[length=2.6mm]}, thick},
    node distance=2.6cm]
  \node[world] (w1) {$w_1$};
  \node[world, right=of w1] (w2) {$w_2$};
  \node[world, right=of w2] (w3) {$w_3$};
  \draw[flecha] (w1) -- (w2);
  \draw[flecha] (w1) to[bend left=28] (w3);
  \draw[flecha] (w2) -- (w3);
  \draw[flecha] (w3) to[loop right, looseness=6] (w3);
  \node[below=2mm of w1] {$p$};
  \node[below=2mm of w2] {$p, q$};
  \node[below=2mm of w3] {$q$};
\end{tikzpicture}
\caption{Un modelo de Kripke con tres mundos; bajo cada mundo se listan las
variables que su valoración hace verdaderas. Se cumple, por ejemplo,
$M, w_1 \vDash \pos p$ (hay un mundo accesible, $w_2$, donde vale $p$);
en cambio $M, w_1 \nvDash \nec p$, porque $p$ falla en el mundo accesible
$w_3$. También $M, w_2 \vDash \nec q$: el único mundo accesible desde $w_2$
es $w_3$, y allí vale $q$. En $w_3$, cuyo único sucesor es él mismo, valen
tanto $\nec q$ como $\pos q$; sin el bucle, $w_3$ no tendría sucesores y
$\nec q$ seguiría siendo verdadera, pero vacuamente, mientras que $\pos q$
pasaría a ser falsa (véase la observación siguiente).}
\label{fig:kripke-ejemplo}
\end{figure}

\begin{observacion}
En un mundo sin sucesores, $\nec\varphi$ es verdadera vacuamente para
cualquier $\varphi$ (¡incluida $\bot$!), pues la condición ``en todo mundo
accesible\dots'' no tiene nada que comprobar. Dualmente, $\pos\varphi$ es
falsa para toda $\varphi$. Este fenómeno, que puede resultar chocante, es
inofensivo en $\K$ y volverá a aparecer en la formalización: el modelo con
el que probaremos la consistencia del sistema es, precisamente, un único
mundo sin sucesores.
\end{observacion}

\begin{definicion}[Validez]\label{def:validez}
Una fórmula $\varphi$ es \emph{válida en un modelo} $M$, escrito
$M \vDash \varphi$, si $M, w \vDash \varphi$ para todo mundo $w$ de $M$;
es \emph{válida en un marco} $F$ si es válida en todo modelo sobre $F$; y es
\emph{válida} (sin más), escrito $\vDash \varphi$, si es válida en todo
modelo de Kripke.
\end{definicion}

Una advertencia terminológica. Blackburn \cite[Def.~1.21]{blackburn} llama
a la primera noción \emph{verdad global} (o \emph{universal}) en el modelo,
y reserva la palabra ``válida'' para los marcos y las clases de marcos
\cite[Def.~1.28]{blackburn}: solo al cuantificar sobre todas las
valoraciones se abstrae de la información contingente que aporta $V$ y se
llega a una noción propiamente lógica. Hemos preferido un único nombre para
los tres niveles, que es también el que usa el código (\lean{ValidInModel},
\lean{ValidInFrame} y \lean{Valid}), pero conviene tener presente la
distinción al consultar el libro.

La noción central para nosotros es la última, porque es la contrapartida
semántica de la demostrabilidad en $\K$. La validez en un marco, que no
depende de la valoración, es la puerta de entrada a las lógicas modales más
fuertes: por ejemplo, $\nec p \limp p$ (con $p$ una variable) es válida en
un marco exactamente cuando este es reflexivo, y añadirla como axioma da la
lógica $\mathbf{T}$.\footnote{De nuevo importa escribir $\nec p \limp p$ y
no $\nec\varphi \limp \varphi$: la instancia $\nec\top \limp \top$ es válida
en todo marco. Lo que caracteriza a los marcos reflexivos es la validez de
la fórmula con una variable o, equivalentemente, la de todas las instancias
del esquema. Estas correspondencias entre fórmulas y propiedades de los
marcos son el tema del capítulo 3 de \cite{blackburn}.} En este trabajo la
definimos por completitud del planteamiento, pero no la explotamos.

Dos ejemplos orientan sobre qué es válido y qué no.

La fórmula
$\nec(\varphi \limp \psi) \limp (\nec\varphi \limp \nec\psi)$ (el axioma K)
es válida: si en todo mundo accesible valen $\varphi \limp \psi$ y
$\varphi$, en todo mundo accesible vale $\psi$; es un ``modus ponens mundo a
mundo''.

En cambio, $\nec p \limp p$ no es válida: en un modelo con dos
mundos $w \mathrel{R} v$, sin más flechas, con $p$ verdadera solo en $v$, se
tiene $M, w \vDash \nec p$ pero no $M, w \vDash p$. Como $\K$ demostrará
exactamente las fórmulas válidas, este contraejemplo muestra también que
$\nec p \limp p$ no es un teorema de $\K$.

\section{Corrección y completitud: el mapa de la demostración}
\label{sec:mapa-completitud}

Tenemos dos nociones de ``verdad lógica'': la sintáctica
($\proves \varphi$) y la semántica ($\vDash \varphi$). Los dos teoremas
centrales de este trabajo afirman que coinciden.

\begin{teorema}[Corrección]\label{teo:correccion-mat}
Si $\proves \varphi$, entonces $\vDash \varphi$.
\end{teorema}

\begin{teorema}[Completitud]\label{teo:completitud-mat}
Si $\vDash \varphi$, entonces $\proves \varphi$.
\end{teorema}

Son los enunciados que formalizaremos (\lean{soundness} y
\lean{completeness} en el capítulo \ref{sec:formalizacion}). Conviene saber
que Blackburn demuestra algo más fuerte \cite[Teorema~4.23]{blackburn}, la
llamada completitud \emph{fuerte}: para todo conjunto de fórmulas $\Sigma$,
si $\varphi$ es consecuencia semántica de $\Sigma$ (es decir, $\varphi$ es
verdadera en todo mundo de todo modelo en el que sean verdaderas todas las
fórmulas de $\Sigma$), entonces $\varphi$ se deduce de $\Sigma$ en $\K$. El
teorema \ref{teo:completitud-mat} es el caso $\Sigma = \emptyset$. En este
trabajo nos limitamos a la versión débil; la construcción que describimos a
continuación es, en cualquier caso, la misma que da la fuerte.

La corrección es la dirección fácil, y su demostración es una inducción
sobre la derivación de $\proves \varphi$: basta comprobar que los axiomas
(las tautologías y el axioma K) son válidos y que las reglas (MP y N)
preservan la validez. El único punto delicado es el caso (CPL): hay que
conectar la noción booleana de tautología (donde las subfórmulas $\nec\psi$
son átomos opacos) con la satisfacción en un mundo concreto. La clave es que,
fijado un mundo $w$, podemos ``leer'' en él una valoración booleana ---la
que asigna a cada átomo y a cada $\nec\psi$ su valor de verdad en $w$--- y
comprobar, por inducción estructural, que evaluar clásicamente con esa
valoración coincide con la satisfacción en $w$. Este lema puente
(\lean{cpl_eval_iff_satisfies} en la formalización) hace todo el trabajo.

Una consecuencia inmediata de la corrección es la \emph{consistencia} del
sistema: $\nvdash \bot$, pues $\bot$ falla en cualquier mundo de cualquier
modelo (y existe al menos un modelo con un mundo).

La completitud es sustancialmente más profunda: hay que fabricar, para cada
fórmula no demostrable, un modelo que la refute. La estrategia clásica, que
seguimos de \cite{blackburn}, construye de una vez un \emph{modelo
canónico} que refuta simultáneamente todas las fórmulas no demostrables. El
plan, cuyos pasos formalizaremos uno a uno en la sección
\ref{sec:form-completitud}, es el siguiente:

\begin{enumerate}
  \item \textbf{Deducibilidad desde hipótesis.} Se generaliza
        $\proves \varphi$ a una relación $\Gamma \proves \varphi$
        (``$\varphi$ se deduce del conjunto de hipótesis $\Gamma$''),
        con su \emph{teorema de deducción}: si
        $\Gamma \cup \{\varphi\} \proves \psi$ entonces
        $\Gamma \proves \varphi \limp \psi$. Aquí nos separamos de
        \cite[Def.~4.4]{blackburn}, que define $\Gamma \proves \varphi$
        mediante conjunciones finitas de premisas: o bien
        $\proves \varphi$ directamente, o bien
        $\proves (\psi_1 \land \dots \land \psi_n) \limp \varphi$ para
        ciertas $\psi_1, \dots, \psi_n \in \Gamma$; nosotros la definimos
        inductivamente, con tres reglas (usar un teorema, usar una
        hipótesis, modus ponens). Ambas definiciones son equivalentes,
        pero la inductiva es mucho más cómoda de formalizar: evita
        manipular listas de fórmulas y conjunciones anidadas, y regala un
        principio de inducción sobre deducciones. Es importante notar que
        la regla de necesitación no se incluye entre las reglas de
        deducción desde hipótesis: de $\Gamma \proves \varphi$ no debe
        seguirse $\Gamma \proves \nec\varphi$ (mis hipótesis pueden ser
        verdades contingentes). Es exactamente la trampa de la observación
        1.41 de \cite{blackburn} que comentamos en la sección
        \ref{sec:sistemaK}. En la definición por conjunciones finitas el
        peligro no existe, porque (N) solo se aplica al teorema
        $\proves (\psi_1 \land \dots \land \psi_n) \limp \varphi$; en la
        inductiva hay que excluirla a propósito.
  \item \textbf{Consistencia y conjuntos maximales.} Un conjunto $\Gamma$ es
        \emph{consistente} si $\Gamma \nvdash \bot$, y \emph{maximal
        consistente} (MCS) si además ninguna extensión propia suya es
        consistente \cite[Def.~4.15]{blackburn}. Los MCS se comportan como
        ``descripciones completas de un mundo'' \cite[Prop.~4.16]{blackburn}:
        están deductivamente cerrados, para cada $\varphi$ contienen
        exactamente una de $\varphi$, $\lnot\varphi$, y contienen
        $\varphi \limp \psi$ si y solo si contener $\varphi$ los obliga a
        contener $\psi$.
  \item \textbf{Lema de Lindenbaum.} Todo conjunto consistente se extiende a
        un MCS. La prueba de \cite[Lema~4.17]{blackburn} enumera las fórmulas
        ($\varphi_0, \varphi_1, \dots$) y las va añadiendo una a una cuando
        no rompen la consistencia; esto exige que el lenguaje sea numerable.
        Nosotros seguimos el camino alternativo del lema de Zorn: la familia
        de los conjuntos consistentes, ordenada por inclusión, tiene la
        propiedad de que toda cadena está acotada (por su unión, que es
        consistente porque toda deducción usa finitas hipótesis), luego
        posee elementos maximales sobre cualquier consistente dado. Se
        elimina así la hipótesis de numerabilidad, y la formalización se
        beneficia del lema de Zorn ya disponible en Mathlib.
  \item \textbf{El modelo canónico.} Sus mundos son todos los MCS;
        la accesibilidad declara $\Gamma \mathrel{R} \Delta$ cuando $\Delta$
        cumple todas las ``obligaciones'' modales de $\Gamma$ (si
        $\nec\psi \in \Gamma$ entonces $\psi \in \Delta$); y la valoración
        hace verdadera $p$ en $\Gamma$ exactamente cuando $p \in \Gamma$.
        Blackburn \cite[Def.~4.18]{blackburn} define la accesibilidad
        canónica con $\pos$ (si $\psi \in \Delta$ entonces
        $\pos\psi \in \Gamma$) y demuestra en su lema 4.19 que equivale a
        la formulación con $\nec$ que adoptamos aquí, la natural cuando
        $\nec$ es el primitivo.
  \item \textbf{Lema de existencia.} Si $\nec\varphi \notin \Gamma$ (con
        $\Gamma$ un MCS), existe un MCS $\Delta$ accesible desde $\Gamma$
        con $\varphi \notin \Delta$. Es la forma dual del lema 4.20 de
        \cite{blackburn} (allí: si $\pos\varphi \in \Gamma$, existe
        $\Delta$ accesible con $\varphi \in \Delta$), y la demostración es
        la misma. Se construye aplicando Lindenbaum al
        conjunto $\{\psi : \nec\psi \in \Gamma\} \cup \{\lnot\varphi\}$,
        cuya consistencia hay que comprobar; el ingrediente clave es que si
        $\{\psi : \nec\psi \in \Gamma\} \proves \varphi$ entonces
        $\Gamma \proves \nec\varphi$ (todas las reglas de la deducción se
        pueden ``encajonar'' gracias al axioma K y a la necesitación).
  \item \textbf{Lema de verdad.} En el modelo canónico, verdad y pertenencia
        coinciden: $M^c, \Gamma \vDash \varphi \iff \varphi \in \Gamma$,
        para toda fórmula $\varphi$ \cite[Lema~4.21]{blackburn}. Se
        demuestra por inducción estructural; los casos proposicionales
        salen de las propiedades de los MCS, y el caso $\nec\varphi$
        combina la definición de la accesibilidad canónica (una dirección)
        con el lema de existencia (la otra).
  \item \textbf{Conclusión, por contraposición.} Si $\nvdash \varphi$,
        entonces $\{\lnot\varphi\}$ es consistente, Lindenbaum lo extiende a
        un MCS $\Gamma$, y el lema de verdad da
        $M^c, \Gamma \nvDash \varphi$: la fórmula $\varphi$ no es válida.
        Es el argumento de los teoremas 4.22 y 4.23 de \cite{blackburn},
        reducido al caso de una sola fórmula.
\end{enumerate}

AAA
